% !TeX root = thesis.tex
\chapter{Conclusions and Recommendations}\label{conclusion}
\section{Conclusions}
For the cohort studied, an efficient transfer depended on which measurements the receiver required, whether a payload already existed, and how an interrupted upload would be completed. Profiling showed why the first decision matters. A missingness filter could remove more than half the predictors while saving only about a third of the rewritten CSV. Lossless compression preserved the values in the columns sent. PCA and the autoencoder offered smaller representations at some settings, but introduced reconstruction error and, particularly for the autoencoder, a fitting cost.

All 82 feasible package configurations passed the receiver recovery checks using SCP. The paired follow-up experiments compared the time each transport took to deliver the same prepared file on the KNUST route. For direct files, the median within-block time to verified delivery favoured rsync for the source CSV, gzip, Parquet and Zstandard, and SCP for the rewritten CSV (Table~\ref{tab:phase2-paired}). For uninterrupted prepared archives, SCP had the shorter paired time for Parquet and gzip, by 1.09~s and 0.46~s respectively. Rsync had a 1.06~s advantage for the larger AE-16 archive (Table~\ref{tab:paired-transport}).

The clearest transport choice arose when the connection broke. Both methods recovered the 42.91~MB AE-16 archive in every interrupted trial. Rsync reused the partial copy at the receiver and reached verified publication sooner in all four matched pairs, with a median saving of 10.46~s. For a transfer of this kind on the KNUST route, rsync is the preferred method. For a single uninterrupted delivery of either of the two smaller prepared archives tested, the measurements favour SCP. These conclusions concern the time until the verified file was available at the receiver. Unpacking and reconstruction followed that endpoint.

\section{Recommendations for research data transfers}
\begin{enumerate}
\item Where the receiving team requires the full set of recorded values, begin with a lossless representation. Parquet is a useful first benchmark. Compare it with gzip on the intended route and include the time required to produce a usable file at the receiver.
\item Before removing a variable, agree with the receiving analyst which observations may be omitted. Record excluded columns, the missingness rule and any protected variables in the manifest. A checksum can establish that the selected file arrived intact, but it cannot recover an omitted measurement.
\item Where approximate reconstruction is acceptable, compare PCA and autoencoders at the same feasible dimension. Examine errors on observed values as well as neighbourhood overlap and the exploratory range screens. Charge fitting and encoding when the payload is created for that delivery.
\item Use SCP for a one-off transfer of the smaller prepared Parquet or gzip archive on a stable link like the one tested. Where a large upload may be interrupted, use rsync over SSH and retain the partial file between attempts. In both cases, verify the archive against its SHA-256 sidecar before releasing it to the receiving team.
\end{enumerate}

\section{Further work}
A useful extension would repeat the paired transfers on routes from participating hospitals, at different times of day and with several interruption points. Varying the patient count while holding the representation and transport fixed would test how delivery cost grows with cohort size. A separate clinical study could compare the original and reconstructed data on a prespecified patient-level task, with patient-level separation and clinical review. Direct feature aggregation, repeated use of a trained encoder and integration with a receiving institution's routine data systems also warrant study.
